\begin{abstract}
LLM Agents route each task through model, tool and Web services, and their encrypted traffic leaks queries, task types and identities. Defenses for this traffic are emerging, but they are judged with website-fingerprinting methods that assume one link is the unit of protection and that cover traffic has no side effects. We systematize fifteen Agent studies (forty-four settings) against twenty-two traffic-analysis anchors and show that neither assumption holds. Tool feedback re-enters model requests, so a task is exposed on several independently controlled links. Agent cover is executed as application work, so its privacy result is incomplete without the task's outcome; no mapped privacy setting reports both an adapted attacker and matched utility. Executed cover also has a timing signature. Cover inserted outside the model loop lacks the planner's inference latency. On 1,440 held-out MCP executions from two LLM planners, a label-free gap rule identifies every decoy. Removing the decoys restores task recognition from 61--82 to 107--116 of 120 tasks, and the attack withstands 200\,ms of timestamp jitter. Dispatching decoys on the planner's own timing distribution removes this advantage and returns the best attack to the size-only level, at a cost of 12--16\% added latency. Cover semantics, by contrast, move size-only recognition by at most ten tasks but decide business-state integrity. We contribute access-aware attack and defense maps, a classical-to-Agent transfer matrix, and an evaluation procedure that binds residual inference to task, state and latency outcomes.
\end{abstract}
